\chapter{Discussion}
\label{ch:discussion}

\section{Why no decision is certified}
\label{sec:disc-ltt}

The triage layer works as designed, and its statistical code reproduces the worked example of the
design exactly. What it lacks is a score that ranks the nodules well enough. Two quantities decide
whether an auto-benign group can be certified at 5\%:

\begin{itemize}
  \item \textbf{Purity:} how few cancers fall among the lowest-scored nodules. This depends on the
        model.
  \item \textbf{Size:} how many nodules that group contains in the calibration set. With 524
        calibration nodules, a 15\% tail has 78 nodules; LTT needs 93 if one of them is
        malignant and 59 if none is.
\end{itemize}

The best language model (run E) has the purity in its lowest 78 nodules but not the size. XGBoost,
with 2 to 11 cancers among its 78 lowest-scored nodules, is at best in the same position. So there are two levers: a model whose low-score tail is cleaner, and a
larger calibration set. Lowering the requirement is a third: the miss-rate variant and Mondrian
split conformal prediction show that a guarantee of a different kind, on the share of cancers
missed rather than on each benign call, is already within reach for about a fifth of the nodules.

The errors that block the guarantee are stable. Most wrong benign calls stay wrong under noise
and under feature removal, so asking the model again does not reveal them. One possible reason is
in the labels: they come from biopsy, and the authors of the dataset report that some biopsy
results disagree with the ultrasound appearance~\citep{gong2022acl}. A malignant nodule that looks
benign on ultrasound cannot be recognised from image features alone, by any model. This is a
hypothesis; it has not been tested here.

\section{What the language model adds}
\label{sec:disc-llm}

The language model learns from the prompts, and it improves when the features improve. It also
appears to use the size of the numbers, not only their order (Section~\ref{sec:results-format}). But with the
same features it stays about 0.05 AUC behind XGBoost, and so far nothing shows that it uses the
\emph{names} of the features, which is the motivation of Ra et al.

Ra et al.\ report that their language model beats all classical baselines. Several points of their
evaluation make this comparison weaker than it looks:

\begin{itemize}
  \item \textbf{Feature names.} Replacing the names with generic labels lowers their internal test
        AUC from 0.700 to 0.695 and their external test AUC from 0.862 to 0.849. These differences
        are called significant, but no test or interval is given.
  \item \textbf{AUC of the baselines.} For all seven scikit-learn baselines, the reported AUC is
        within 0.004 of the accuracy (for example 0.631 and 0.631 for the random forest), while
        for their own model and for TransTab the AUC is clearly higher than the accuracy (0.700
        against 0.671). This pattern is expected if the baselines' AUC was computed from 0/1
        predictions, which turns AUC into balanced accuracy. The comparison would then favour the
        language model.
  \item \textbf{Fine-tuning.} Training only the classification head of the frozen model gives an
        AUC of 0.699, against 0.700 with LoRA; without mRMR, head-only training gives 0.710, the
        highest AUC in their ablation table.
  \item \textbf{Labels.} Their labels are BI-RADS categories (3 benign, 4--6 malignant), that is a
        radiologist's assessment of the image, not a biopsy result.
  \item \textbf{Size and uncertainty.} The internal test set has 235 images and the external one
        about 55 after half of INbreast was used for retraining. No intervals are reported, and
        the model was trained for over 500 epochs and selected on validation.
\end{itemize}

The results of this thesis are in line with a cautious reading of theirs: a language model can
classify radiomics prompts at about the level of classical models, and a clear advantage remains
to be shown. Two steps can test the claim directly: censoring the feature names, and giving the
model plain-language meanings of the features (Chapter~\ref{ch:future}).

\section{Limitations}
\label{sec:disc-limitations}

\begin{itemize}
  \item \textbf{One 2D image per nodule.} A nodule is a 3D structure seen in several planes and in
        real time; one still image gives a partial view.
  \item \textbf{Ultrasound images.} Brightness depends on the device, its settings and the
        operator. Per-image rescaling removes part of this variation, and with it part of the
        information on absolute echogenicity.
  \item \textbf{Clinical signs not captured.} Several TI-RADS signs~\citep{tessler2017} are not
        measured by the features: echogenic foci (microcalcifications), margin details, a
        taller-than-wide shape in the transverse plane, and echogenicity relative to the
        surrounding thyroid tissue.
  \item \textbf{Patient overlap.} The pooled split is drawn at image level, so the absolute test
        metrics are likely optimistic for new patients (Section~\ref{sec:data-patients}).
  \item \textbf{Single runs.} Each setting was trained once, with one seed. Differences below
        about 0.02 AUC are within seed-to-seed variation.
  \item \textbf{Hyperparameters.} The learning rate was chosen for the 1.5B model and reused for
        7B; XGBoost was not tuned.
  \item \textbf{One dataset.} All results come from TN3K.
  \item \textbf{Exploratory analyses.} The miss-rate variant was added after the main results.
\end{itemize}
